% Chemotaxis compressible Navier–Stokes (vascular network formation)
% Converted from notes/pde/case-chemotaxis.html. Edit freely: this file is now the master copy.

\chapter{Chemotaxis compressible Navier--Stokes equations}\label{p:case-chemotaxis}

\begin{byline}Based on T. Huang, B. Rababah, A. Rawashdeh, \emph{On existence of chemotaxis compressible Navier--Stokes equations modeling vascular network formation}, J. Differential Equations 480 (2026) 114655. \href{https://doi.org/10.1016/j.jde.2026.114655}{doi:10.1016/j.jde.2026.114655}\end{byline}

This page walks through a complete existence theory for one nonlinear system. The aim is to show how the standard tools of PDE analysis fit together in a real argument. Every tool is linked to its own concept page, which states it precisely and points to a proof. Steps are tagged \routine{} when a standard tool is used in the standard way, and \newidea{} when the argument does something that cannot be taken off the shelf.

\begin{conceptsbox}
\begin{multicols}{2}\small\raggedright
\begin{itemize}[leftmargin=*,itemsep=1pt]
\item Weak derivatives and Sobolev spaces~(\S\ref{p:sobolev-spaces})
\item Spaces of time-dependent functions~(\S\ref{p:bochner-spaces})
\item Weak and weak* convergence~(\S\ref{p:weak-convergence})
\item Hölder, Young, and L\textsuperscript{\textit{p}} interpolation~(\S\ref{p:elementary-inequalities})
\item Sobolev embeddings~(\S\ref{p:sobolev-embedding})
\item Gagliardo--Nirenberg inequalities~(\S\ref{p:gagliardo-nirenberg})
\item Gronwall's inequality and continuity arguments~(\S\ref{p:gronwall})
\item Elliptic regularity: the Lamé system and the Neumann Laplacian~(\S\ref{p:elliptic-regularity})
\item Maximal regularity for parabolic equations~(\S\ref{p:parabolic-regularity})
\item The transport (continuity) equation~(\S\ref{p:transport-equation})
\item Strong, weak, and renormalized solutions~(\S\ref{p:notions-of-solution})
\item Energy estimates and a priori bounds~(\S\ref{p:energy-estimates})
\item The Faedo--Galerkin method~(\S\ref{p:galerkin})
\item Uniqueness by difference estimates~(\S\ref{p:uniqueness-energy})
\item Compactness: Rellich--Kondrachov and Aubin--Lions~(\S\ref{p:compactness})
\item Regularized approximation schemes~(\S\ref{p:regularized-approximation})
\item Lions--Feireisl theory for compressible flow~(\S\ref{p:lions-feireisl})
\end{itemize}
\end{multicols}
\end{conceptsbox}

\section{The model}

Endothelial cells, spread on a substrate, organize themselves into a network of blood vessels. They move like a compressible fluid and are attracted up the gradient of a chemical that they themselves secrete. In a bounded domain $\Omega \subset \R^3$ with smooth boundary, the unknowns are the cell density $\rho \ge 0$, the cell velocity $u \in \R^3$, and the chemoattractant concentration $\phi \ge 0$:

\begin{equation}\label{eq:system}
  \left\{
  \begin{aligned}
    &\rho_t + \dv(\rho u) = 0, \\
    &(\rho u)_t + \dv(\rho u \otimes u) + \nabla p(\rho) - \mu \Delta u - (\lambda+\mu)\nabla \dv u
       = \chi \rho \nabla \phi - \alpha \rho u, \\
    &\sigma \phi_t = D \Delta \phi + a \rho - b \phi,
  \end{aligned}
  \right.
  \end{equation}

with pressure $p(\rho) = \frac{\theta(\gamma-1)}{\gamma}\rho^\gamma$, $\gamma > 1$. The coefficients satisfy $D, \chi, \alpha, a, \theta > 0$, $\mu > 0$, $3\lambda + 2\mu > 0$, and $b, \sigma \ge 0$.

\begin{itemize}
\item $\chi \rho \nabla \phi$ is the chemotactic force: cells are pushed up the gradient of $\phi$.
\item $\alpha \rho u$ is friction with the substrate (damping).
\item $a\rho$ is secretion of the chemical by the cells; $b\phi$ is its decay.
\item The pressure models the fact that closely packed cells resist compression.
\end{itemize}

The boundary and initial conditions are

\begin{equation}\label{eq:bc}
    u = 0, \quad \frac{\partial \phi}{\partial \nu} = 0 \ \text{ on } \partial\Omega, \qquad
    (\rho, u, \phi)|_{t=0} = (\rho_0, u_0, \phi_0).
  \end{equation}

Setting $\mu = \lambda = 0$ gives the inviscid hyperbolic--parabolic model of Gamba et al. (2003), which reproduces the network patterns seen in experiments. Here we keep the viscosity.

\subsection*{Two difficulties to keep in mind}

\textbf{Vacuum.} The density may vanish on an open set. Where $\rho = 0$ the momentum equation degenerates: the term $\rho u_t$ gives no control of $u_t$. Both theorems below allow this.

\textbf{An energy with no sign.} The natural energy contains the cross term $-\chi\rho\phi$, which can be negative and large. Before any estimate can be used, one must show that the energy is bounded below.

\section{The main results}

\begin{thmbox}{Theorem A (local strong solutions)}
Let $q \in (3,6]$. Suppose $\rho_0 \ge 0$, $\rho_0 \in W^{1,q}$, $u_0 \in D^2 \cap D^1_0$, $0 \le \phi_0 \in H^2$, and the compatibility condition

\begin{equation}\label{eq:compat}
      \mu \Delta u_0 + (\lambda+\mu)\nabla \dv u_0 - \nabla p(\rho_0) + \chi \rho_0 \nabla \phi_0 - \alpha \rho_0 u_0
      = \sqrt{\rho_0}\, g \quad \text{for some } g \in L^2(\Omega)
    \end{equation}

holds. Then there are $T_0 > 0$ and a unique strong solution of \eqref{eq:system}--\eqref{eq:bc} on $\Omega \times (0, T_0]$.
\end{thmbox}

\begin{proofslot}Proof of Theorem A: to be written. The map in Section 4 lists the steps in order.\end{proofslot}

\begin{thmbox}{Theorem B (global weak solutions)}
Let $\gamma > 3/2$. Suppose $0 \le \rho_0 \in L^\gamma$ with $\rho_0 \not\equiv 0$, $\sqrt{\rho_0}\, u_0 \in L^2$, and $0 \le \phi_0 \in W^{2 - 2/\gamma,\, \gamma}$. Then for every $T > 0$ there is a finite-energy weak solution of \eqref{eq:system}--\eqref{eq:bc} on $\Omega \times (0,T)$.
\end{thmbox}

\begin{proofslot}Proof of Theorem B: to be written. The map in Section 5 lists the steps in order.\end{proofslot}

What "strong solution" and "finite-energy weak solution" mean, with the exact function spaces, is on Strong, weak, and renormalized solutions~(\S\ref{p:notions-of-solution}); the spaces themselves are on Sobolev spaces~(\S\ref{p:sobolev-spaces}) and spaces of time-dependent functions~(\S\ref{p:bochner-spaces}).

\section{The energy structure}

Both theorems rest on one identity. Define

\begin{equation}\label{eq:energy}
    E(t) = \int_\Omega \Big( \tfrac12 \rho |u|^2 + \tfrac{\theta}{\gamma}\rho^\gamma
      + \tfrac{\chi D}{2a}|\nabla\phi|^2 + \tfrac{\chi b}{2a}\phi^2 - \chi\rho\phi \Big)\, dx .
  \end{equation}

\begin{step}{\routine{}The energy identity (Lemma 2.1)}
For smooth solutions, $\displaystyle \frac{d}{dt}E(t) = -\int_\Omega \Big( \tfrac{\chi\sigma}{a}|\phi_t|^2 + \alpha\rho|u|^2 + \mu|\nabla u|^2 + (\lambda+\mu)|\dv u|^2 \Big)dx \le 0.$

Test the momentum equation with $u$ and the $\phi$-equation with $\frac{\chi}{a}\phi_t$. The chemotactic work $\int \chi\rho\nabla\phi\cdot u$ is rewritten with the continuity equation as $\frac{d}{dt}\int\chi\rho\phi - \int\chi\rho\phi_t$, and the second piece cancels against the $\phi$-equation. Choosing the multipliers so that the cross terms cancel is the whole point.

\uses{energy estimates~(\S\ref{p:energy-estimates}).}
\end{step}

\begin{stepnew}{\newidea{}The energy is bounded below for $\gamma > 4/3$ (Proposition 2.2)}
Let $E_+$ be the energy without the cross term. Using only conservation of mass, $\int\rho = m$, one shows $\abs{\int\chi\rho\phi} \le C m^{q_0} + \frac12 E_+(t)$, so $E \ge \frac12 E_+ - Cm^{q_0}$. The cross term is split by Hölder, then $\norm{\rho}_{L^l}$ is interpolated between the conserved $L^1$ norm and the energy's $L^\gamma$ norm, and $\norm{\phi}_{L^{l'}}$ between $L^2$ and $H^1$. The exponents close (the power of $E_+$ stays below 1) exactly when $\gamma > 4/3$. Earlier work needed $\gamma > 8/5$. There are three cases, $\gamma > 2$, $\gamma = 2$, and $4/3 < \gamma < 2$; the last is where the exponent bookkeeping is delicate.

\uses{Hölder, Young, and L\textsuperscript{\textit{p}} interpolation~(\S\ref{p:elementary-inequalities}), Gagliardo--Nirenberg~(\S\ref{p:gagliardo-nirenberg}), Sobolev embeddings~(\S\ref{p:sobolev-embedding}).}
\end{stepnew}

\section{Map of the proof of Theorem A}

The strategy: first prove bounds that every smooth solution with $\rho_0 \ge \delta > 0$ must satisfy, with constants independent of $\delta$. Then build approximate solutions that satisfy the same bounds and pass to the limit. All the higher-order bounds are collected in one quantity,

\[
    \Phi(t) = \sup_{0 \le s \le t}\Big( \norm{\rho(s)}_{H^1 \cap W^{1,q}} + \norm{\nabla u(s)}_{L^2}
      + \norm{\sqrt{\rho}\,u_t(s)}_{L^2} + \norm{\phi(s)}_{H^2} + 1 \Big),
  \]

and the goal of Section 3 of the paper is the inequality $\Phi(t) \le \exp\big[C\mathcal{M} + C\int_0^t M(\Phi(s))\,ds\big]$, where $\mathcal{M}$ depends only on the data (through $\norm{g}_{L^2}$) and $M$ is some increasing function. Such an inequality forces $\Phi$ to stay bounded for a short time.

\begin{step}{\routine{}Step 1. Second derivatives of $u$ from the equation (Lemma 3.2)}
Read the momentum equation as $\Lame u = \rho u_t + \rho u\cdot\nabla u + \nabla p(\rho) - \chi\rho\nabla\phi + \alpha\rho u$ for the Lamé operator $\Lame$. Elliptic regularity bounds $\norm{\nabla u}_{H^1}$ by the right-hand side in $L^2$. The convection term produces $\norm{\nabla u}_{H^1}$ again, at a power below 2, so it is absorbed. Result: $\norm{\nabla u}_{H^1} \le M(\Phi)$.

\uses{elliptic regularity~(\S\ref{p:elliptic-regularity}), Sobolev embeddings~(\S\ref{p:sobolev-embedding}), Gagliardo--Nirenberg~(\S\ref{p:gagliardo-nirenberg}), Young's inequality~(\S\ref{p:elementary-inequalities}).}
\end{step}

\begin{step}{\routine{}Step 2. The time derivative $\sqrt\rho\, u_t$ (Lemmas 3.3--3.4)}
Differentiate the momentum equation in $t$, test with $u_t$, and estimate the six resulting terms. The value at $t = 0$, $\norm{\sqrt{\rho_0}\,u_t(0)}_{L^2}$, is exactly $\norm{g}_{L^2}$ from the compatibility condition \eqref{eq:compat}. This is why that condition is needed: without it the estimate has no starting value when $\rho_0$ vanishes somewhere.

\uses{energy estimates~(\S\ref{p:energy-estimates}), compatibility conditions~(\S\ref{p:notions-of-solution}), Sobolev embeddings~(\S\ref{p:sobolev-embedding}).}
\end{step}

\begin{step}{\routine{}Step 3. The density (Lemma 3.5)}
The continuity equation is a transport equation: it carries $\rho$ along the flow and never smooths it. Its $W^{1,q}$ norm grows at most like $\exp\big(C\int_0^t \norm{\nabla u}_{H^1 \cap D^{1,q}}\big)$. So one needs $\nabla u \in L^1(0,T; W^{1,q})$ with $q > 3$, which comes from the $W^{2,q}$ estimate for the Lamé system and Step 2.

\uses{the transport equation~(\S\ref{p:transport-equation}), elliptic regularity~(\S\ref{p:elliptic-regularity}).}
\end{step}

\begin{step}{\routine{}Step 4. The chemical (Lemma 3.6)}
Take the gradient of the $\phi$-equation, test with $\nabla\phi_t$, and use the Neumann condition. This bounds $\norm{\Delta\phi}_{L^2}$, hence $\norm{\phi}_{H^2}$ by elliptic regularity for the Neumann Laplacian.

\uses{parabolic regularity~(\S\ref{p:parabolic-regularity}), elliptic regularity~(\S\ref{p:elliptic-regularity}).}
\end{step}

\begin{step}{\routine{}Step 5. Closing the estimate}
Combine Steps 1--4 into the inequality for $\Phi$ above. A Gronwall-type continuity argument gives $T_0 > 0$, depending only on the data, with $\sup_{[0,T_0]}\Phi \le C\exp(C\norm{g}_{L^2})$.

\uses{Gronwall and continuity arguments~(\S\ref{p:gronwall}).}
\end{step}

\begin{step}{\routine{}Step 6. Construction (Section 4.1)}
Lift the density to $\rho_0 + \delta$ so there is no vacuum, adjusting $u_0$ so that the compatibility condition still holds with the same $g$. Project $u$ onto finitely many eigenfunctions of the Lamé operator; solve the continuity and $\phi$ equations exactly and the projected momentum equation as an ODE system (the mass matrix is invertible because $\rho \ge \delta > 0$). The bounds of Step 5 hold uniformly in the number of modes and in $\delta$, so one can pass to weak limits, first in the number of modes and then $\delta \to 0$.

\uses{the Faedo--Galerkin method~(\S\ref{p:galerkin}), weak and weak* convergence~(\S\ref{p:weak-convergence}), spaces of time-dependent functions~(\S\ref{p:bochner-spaces}).}
\end{step}

\begin{step}{\routine{}Step 7. Uniqueness (Section 4.2)}
Subtract two solutions. Estimate the difference in the lower norm $\norm{\sqrt{\rho_1}\,u}_{L^2}^2 + \norm{\rho}_{L^2}^2 + \norm{\phi}_{H^1}^2$, using the regularity of the two solutions only as coefficients. The result is a linear Gronwall inequality with zero initial value, so the difference vanishes.

\uses{uniqueness by difference estimates~(\S\ref{p:uniqueness-energy}), Gronwall's inequality~(\S\ref{p:gronwall}).}
\end{step}

\section{Map of the proof of Theorem B}

Weak solutions are built through a chain of regularized problems, each removing one difficulty, then the regularizations are removed one at a time. The paper concentrates on $3/2 < \gamma < 2$; larger $\gamma$ was covered by earlier work.

\begin{stepnew}{\newidea{}Step 1. The regularized system}
\[
    \left\{
    \begin{aligned}
      &\rho_t + \dv(\rho u) = \varepsilon\Delta\rho, \\
      &(\rho u)_t + \dv(\rho u\otimes u) + \nabla p(\rho) + \delta\nabla\rho^\beta + \varepsilon\nabla\rho\cdot\nabla u \\
      &\qquad = \mu\Delta u + (\lambda+\mu)\nabla\dv u + \chi\rho\nabla\phi - \alpha\rho u, \\
      &\sigma\phi_t + \varepsilon\Delta\big(|\Delta\phi|^{p_0}\Delta\phi\big) = D\Delta\phi + a\rho - b\phi,
      \qquad p_0 = \tfrac{2-\gamma}{\gamma-1}.
    \end{aligned}
    \right.
    \]

The first three regularizations are the standard Feireisl--Novotný--Petzeltová scheme: artificial viscosity $\varepsilon\Delta\rho$, a balancing term $\varepsilon\nabla\rho\cdot\nabla u$, and artificial pressure $\delta\rho^\beta$. The new difficulty: the artificial viscosity creates the term $-\varepsilon\int\phi\,\Delta\rho = -\varepsilon\int\rho\,\Delta\phi$ in the energy balance. Handled in the usual way, it leads to a term like $\varepsilon\int\rho^2$, which the energy controls only when $\gamma \ge 2$. The quasi-linear biharmonic term is designed for the range $3/2 < \gamma < 2$. It adds $\frac{\varepsilon\chi}{a(p_0+2)}\int|\Delta\phi|^{p_0+2}$ to the energy, and $p_0$ is chosen so that $p_0 + 2$ and $\gamma$ are conjugate exponents, so Young's inequality splits $\varepsilon\int\rho\,\Delta\phi$ between $\int\rho^\gamma$ and the new term.

\uses{regularized approximation schemes~(\S\ref{p:regularized-approximation}), Young's inequality~(\S\ref{p:elementary-inequalities}).}
\end{stepnew}

\begin{step}{\routine{}Step 2. Uniform energy bounds (Lemma 5.1)}
The energy identity of Section 3, redone for the regularized system, with the extra terms controlled as in Step 1 and the lower bound of Proposition 2.2. The constants do not depend on the Galerkin index, $\varepsilon$, or $\delta$.

\uses{energy estimates~(\S\ref{p:energy-estimates}), Gronwall's inequality~(\S\ref{p:gronwall}).}
\end{step}

\begin{step}{\routine{}Step 3. Removing the regularizations}
The Galerkin limit and $\varepsilon \to 0$ follow the compressible Navier--Stokes theory, using compactness to get strong convergence where the nonlinear terms need it. The biharmonic term disappears because the energy bounds $\varepsilon\int|\Delta\phi|^{p_0+2}$, so the term itself is $O(\varepsilon^{1/(p_0+2)})$ in the weak sense. The limit $\delta \to 0$ needs a little more integrability of $\rho$ than $L^\gamma$; this is where the pressure and the effective viscous flux enter.

\uses{compactness~(\S\ref{p:compactness}), Lions--Feireisl theory~(\S\ref{p:lions-feireisl}), weak convergence~(\S\ref{p:weak-convergence}).}
\end{step}

\begin{step}{\routine{}Step 4. Why $\gamma > 3/2$}
For $\rho\nabla\phi$ to make sense in the weak formulation, $\nabla\phi$ must lie in $L^{\gamma'}$, $\gamma' = \frac{\gamma}{\gamma-1}$, since $\rho$ is only in $L^\gamma$. Maximal regularity for the $\phi$-equation, with $a\rho \in L^\gamma$ as the source, gives $\phi \in W^{2,1}_\gamma$, so $\nabla\phi \in W^{1,\gamma}$ in space. The Sobolev embedding $W^{1,\gamma} \subset L^{3\gamma/(3-\gamma)}$ is enough when

\[
      \frac{3\gamma}{3-\gamma} > \frac{\gamma}{\gamma-1} \iff 3(\gamma - 1) > 3 - \gamma \iff \gamma > \tfrac32 .
    \]

The same threshold is where the Feireisl theory for compressible Navier--Stokes stops, so both obstructions meet at $3/2$.

\uses{maximal parabolic regularity~(\S\ref{p:parabolic-regularity}), Sobolev embeddings~(\S\ref{p:sobolev-embedding}).}
\end{step}

\section{Exercises}

These rebuild the paper's key computations by hand.

\begin{exercise}
  Derive the energy identity of Lemma 2.1 for smooth solutions: test the momentum equation with $u$ and the $\phi$-equation with $\frac{\chi}{a}\phi_t$, and show how the cross terms cancel. Check that the pressure potential is $\frac\theta\gamma\rho^\gamma$.
  \begin{hint}
    The chemotactic work: $\int\chi\rho\nabla\phi\cdot u = -\int\chi\phi\,\dv(\rho u) = \int\chi\phi\rho_t = \frac{d}{dt}\int\chi\rho\phi - \int\chi\rho\phi_t$. The $\phi$-equation tested with $\frac\chi a\phi_t$ produces $+\int\chi\rho\phi_t$. For the pressure, $P(\rho) = \frac{p(\rho)}{\gamma-1} = \frac\theta\gamma\rho^\gamma$ satisfies $\rho P'(\rho) - P(\rho) = p(\rho)$.
  \end{hint}
  \begin{solution}
    \emph{Momentum.} Test with $u$ ($u = 0$ on $\partial\Omega$). Using the continuity equation, $\int\big((\rho u)_t + \dv(\rho u\otimes u)\big)\cdot u = \frac{d}{dt}\int\frac12\rho|u|^2$. For the pressure, with $P(\rho) = \frac{\theta}{\gamma}\rho^\gamma$ one has $\rho P'(\rho) - P(\rho) = \theta\rho^\gamma - \frac\theta\gamma\rho^\gamma = p(\rho)$, and then $\int\nabla p\cdot u = -\int p\,\dv u = \frac{d}{dt}\int P(\rho)$. The viscous terms give $\mu\norm{\nabla u}^2 + (\lambda+\mu)\norm{\dv u}^2$. So

    \[ \frac{d}{dt}\int\Big(\frac12\rho|u|^2 + \frac\theta\gamma\rho^\gamma\Big) + \mu\norm{\nabla u}^2 + (\lambda+\mu)\norm{\dv u}^2
             = \int\chi\rho\nabla\phi\cdot u - \alpha\int\rho|u|^2. \]

    \emph{Chemotactic work.} $\int\chi\rho\nabla\phi\cdot u = -\int\chi\phi\,\dv(\rho u) = \int\chi\phi\rho_t = \frac{d}{dt}\int\chi\rho\phi - \int\chi\rho\phi_t$.

    \emph{Chemical.} Test $\sigma\phi_t = D\Delta\phi + a\rho - b\phi$ with $\frac\chi a\phi_t$, using $\partial_\nu\phi = 0$:

    \[ \frac{\chi\sigma}{a}\norm{\phi_t}^2 = -\frac{d}{dt}\int\Big(\frac{\chi D}{2a}|\nabla\phi|^2 + \frac{\chi b}{2a}\phi^2\Big) + \int\chi\rho\phi_t. \]

    Add the three: the term $\int\chi\rho\phi_t$ appears with opposite signs and cancels, leaving $\frac{d}{dt}E = -\int\big(\frac{\chi\sigma}{a}|\phi_t|^2 + \alpha\rho|u|^2 + \mu|\nabla u|^2 + (\lambda+\mu)|\dv u|^2\big)$.
  \end{solution}
\end{exercise}

\begin{exercise}
  Proposition 2.2 in its simplest case. Assume $b > 0$ and $\gamma > 2$. Using the mass $m$, show $\big|\int\chi\rho\phi\big| \le \frac12E_+ + Cm$, where $E_+$ contains $\frac\theta\gamma\int\rho^\gamma$ and $\frac{\chi b}{2a}\int\phi^2$.
  \begin{hint}
    $|\int\rho\phi| \le \norm{\rho}_{L^2}\norm{\phi}_{L^2} \le m^\eta\norm{\rho}_{L^\gamma}^{1-\eta}\norm{\phi}_{L^2}$ with $\eta = \frac{\gamma-2}{2(\gamma-1)}$ (Exercise 3 in Hölder, Young, and L\textsuperscript{\textit{p}} interpolation~(\S\ref{p:elementary-inequalities})). Use Young with three exponents: $\norm{\rho}_{L^\gamma}^{1-\eta}$ to the power $\frac{\gamma}{1-\eta}$, $\norm{\phi}_{L^2}$ to the power $2$, and $m^\eta$ to the power $r$ with $\frac{1-\eta}{\gamma} + \frac12 + \frac1r = 1$. You should find $r = \frac1\eta$, which needs $\gamma > 2$.
  \end{hint}
  \begin{solution}
    By Hölder and Exercise 3 in the inequalities section ($\rho \ge 0$, $\int\rho = m$): $\big|\int\rho\phi\big| \le \norm{\rho}_{L^2}\norm{\phi}_{L^2} \le m^\eta\norm{\rho}_{L^\gamma}^{1-\eta}\norm{\phi}_{L^2}$ with $\eta = \frac{\gamma-2}{2(\gamma-1)}$, so $1 - \eta = \frac{\gamma}{2(\gamma-1)}$.

    Young with three factors: $abc \le \frac{a^P}{P} + \frac{b^Q}{Q} + \frac{c^R}{R}$ when $\frac1P + \frac1Q + \frac1R = 1$. Take $a = \varepsilon_1\norm{\rho}_{L^\gamma}^{1-\eta}$ with $P = \frac{\gamma}{1-\eta} = 2(\gamma-1)$ (so $a^P = \varepsilon_1^P\norm{\rho}_{L^\gamma}^\gamma$), $b = \varepsilon_2\norm{\phi}_{L^2}$ with $Q = 2$, and $c = \frac{\chi m^\eta}{\varepsilon_1\varepsilon_2}$. Then $\frac1R = 1 - \frac{1}{2(\gamma-1)} - \frac12 = \frac{\gamma-2}{2(\gamma-1)} = \eta$, positive exactly when $\gamma > 2$, and $c^R = \big(\frac{\chi}{\varepsilon_1\varepsilon_2}\big)^{1/\eta}m$. Hence

    \[ \chi\Big|\int\rho\phi\Big| \le \frac{\varepsilon_1^P}{P}\int\rho^\gamma + \frac{\varepsilon_2^2}{2}\int\phi^2 + C(\varepsilon_1,\varepsilon_2)\,m. \]

    Choose $\varepsilon_1, \varepsilon_2$ with $\frac{\varepsilon_1^P}{P} \le \frac{\theta}{2\gamma}$ and $\frac{\varepsilon_2^2}{2} \le \frac{\chi b}{4a}$ (this uses $b > 0$). The first two terms are then at most half of $\frac\theta\gamma\int\rho^\gamma + \frac{\chi b}{2a}\int\phi^2 \le E_+$, so $\big|\int\chi\rho\phi\big| \le \frac12E_+ + Cm$.
  \end{solution}
\end{exercise}

\begin{exercise}
  The exponent $p_0$ in the regularized system. Show $p_0 + 2 = \frac{\gamma}{\gamma-1}$, the conjugate of $\gamma$. Then use Young to bound $\varepsilon\big|\int\rho\,\Delta\phi\big|$ by $\varepsilon\int\rho^\gamma$ and $\varepsilon\int|\Delta\phi|^{p_0+2}$.
  \begin{hint}
    $p_0 + 2 = \frac{2-\gamma+2\gamma-2}{\gamma-1}$. Then $\rho|\Delta\phi| \le \frac{\rho^\gamma}{\gamma} + \frac{|\Delta\phi|^{\gamma'}}{\gamma'}$.
  \end{hint}
  \begin{solution}
    $p_0 + 2 = \frac{2-\gamma}{\gamma-1} + 2 = \frac{2 - \gamma + 2\gamma - 2}{\gamma-1} = \frac{\gamma}{\gamma-1} = \gamma'$, and $\frac1\gamma + \frac1{\gamma'} = 1$. For $\rho \ge 0$, Young's inequality with exponents $\gamma, \gamma'$ (and a parameter $\kappa > 0$) gives

    \[ \varepsilon\Big|\int\rho\,\Delta\phi\Big| \le \varepsilon\int\rho|\Delta\phi| \le \varepsilon\kappa\int\rho^\gamma + \varepsilon C(\kappa)\int|\Delta\phi|^{p_0+2}. \]

    The regularized energy contains $\frac\theta\gamma\int\rho^\gamma$ and $\frac{\varepsilon\chi}{a(p_0+2)}\int|\Delta\phi|^{p_0+2}$, so both terms on the right are at most a constant times that energy. In the energy balance this gives $\frac{d}{dt}E \le CE$ (plus dissipation), and Gronwall closes the estimate. Without the biharmonic term one would need $\Delta\phi \in L^{\gamma'}$ from somewhere else, and $\gamma' > 2$ when $\gamma < 2$.
  \end{solution}
\end{exercise}

\section{Open questions}

\begin{itemize}
\item \textbf{Vanishing damping.} The construction of weak solutions never uses the damping term $\alpha\rho u$. What happens as $\alpha \to 0$?
\item \textbf{Vanishing viscosity.} Can weak solutions of the inviscid model ($\mu = \lambda = 0$) be obtained as limits of solutions of \eqref{eq:system}?
\item \textbf{The gap $4/3 < \gamma \le 3/2$.} The energy is bounded below there, but global weak solutions are not known to exist.
\item \textbf{Lower dimensions.} The same method should give weak solutions for $\gamma > 4/3$ in two dimensions and $\gamma > 1$ in one dimension.
\end{itemize}

\section*{References}

\begin{enumerate}[label={[\arabic*]}, leftmargin=*]
\item T. Huang, B. Rababah, A. Rawashdeh, On existence of chemotaxis compressible Navier--Stokes equations modeling vascular network formation, \emph{J. Differential Equations} 480 (2026) 114655.
\item X. Huo, A. Jüngel, Global existence and weak-strong uniqueness for chemotaxis compressible Navier--Stokes equations modeling vascular network formation, \emph{J. Math. Fluid Mech.} 26 (2024) 11.
\item Y. Cho, H. J. Choe, H. Kim, Unique solvability of the initial boundary value problems for compressible viscous fluids, \emph{J. Math. Pures Appl.} 83 (2004) 243--275.
\item E. Feireisl, A. Novotný, H. Petzeltová, On the existence of globally defined weak solutions to the Navier--Stokes equations, \emph{J. Math. Fluid Mech.} 3 (2001) 358--392.
\item A. Gamba et al., Percolation, morphogenesis, and Burgers dynamics in blood vessels formation, \emph{Phys. Rev. Lett.} 90 (2003) 118101.
\end{enumerate}
